% s-134-p134-m002; original PDF 134
\Aphorism{s-134-p134-m002}{189}
Nur ein Verhältniss zwischen Männern ist Ruhe, Meeresstille, Sicherheit ohne Verstecke, offenes Gelände ohne die Gefahr peinlicher Überraschungen. Für ein Verhältniss zwischen Mann und Weib gilt der tragische Spruch „Alles\Sperr{ klar}, aber auch Alles\Sperr{ zu Ende}“; die Beziehung zum Weibe muss unerbittlich jenen gefährlichen Punkt passiren, der der Kreuzungspunkt schlimmer Möglichkeiten und Wirrnisse und Enttäuschungen ist. Trennung oder Erniedrigung; ein dritter Weg ist wohl noch denkbar, aber wer ist ihn je gegangen? — Männer, deren Freundschaft gar zu tief, zu ruhevoll geworden ist, suchen schliesslich Bewegung um jeden Preis; darum kommen oft vertrauteste Freunde durch ganz subalterne Anlässe auseinander. Zwischen Mann und Weib werden umgekehrt höchst gebieterische Anlässe, ein Ende zu machen, überhört; aus dieser Stumpfheit, die um Katastrophen herumgeht statt hindurch, macht man gar noch ein Ideal, Treue, und eine Institution, die Ehe.
\par

